\chapter{年报、中报与预告：先分实际性，再讨论空间}

财务更新的第一原则是实际性：先确认数字是否已由正式中报披露，再讨论其持续性。下表中只有恒逸和江波龙为实际H1；其余H1利润范围都不能替代收入、利润率、现金和营运资本的完整证据。

\begin{exhibitbox}[图表4｜主要财务证据，单位：CNY亿元]
\small
\begin{tabularx}{\textwidth}{L{1.1cm}L{1.55cm}R{1.0cm}R{1.4cm}L{1.7cm}X}
\toprule
\textbf{代码} & \textbf{标的} & \textbf{FY2025归母} & \textbf{2026Q1归母/OCF} & \textbf{H1状态} & \textbf{解读} \\
\midrule
000703 & 恒逸石化 & 2.58 & 19.95 / -25.73 & 实际中报：归母59.02，OCF -4.88 [F-000703-H1] & 利润已验证，现金仍须解释。 \\
002379 & 宏桥控股 & 178.64 & 67.58 / 92.38 & 预告：归母150--160 [F-002379-H1] & 方向积极，但收入和H1 OCF未披露。 \\
300014 & 亿纬锂能 & 41.34 & 14.46 / -3.66 & 预告：归母31.30--33.71 [F-300014-H1] & 不能把预告当作完整经营验证。 \\
002497 & 雅化集团 & 6.32 & 3.39 / -4.31 & 预告：归母11--13 [F-002497-H1] & H2利润与现金是全年逻辑的核心。 \\
301308 & 江波龙 & 14.23 & 38.62 / -28.75 & 实际中报：归母105.77，OCF -31.51 [F-301308-H1] & 不能以利润替代现金转换。 \\
\bottomrule
\end{tabularx}
\normalsize
\sourcenote{F-000703-H1：2026-08-11 CNINFO公告1225467439，第10页；F-002379-H1：2026-07-11公告1225418343，第1--2页；F-300014-H1：2026-06-15公告1225371777，第1--2页；F-002497-H1：2026-07-07公告1225411869，第1--2页；F-301308-H1：2026-08-11公告1225467150，第11页。部分可比期间存在重述，未机械引用同比增速。}
\end{exhibitbox}

图表4直接决定图表1的状态：恒逸可验证的是利润而非低位，宏桥、亿纬、雅化只能验证预告方向而不能建立完整分母，江波龙则提醒“利润高”不等于“现金已兑现”。

\section{实际中报只有两份}

本期13只样本中，恒逸与江波龙披露了实际H1报告。其余九只为发行人预告范围，辽宁成大和中国船舶在官方检索中没有找到H1报告或标题匹配的预告。预告可用于确认方向，但不能提供完整的收入、利润率、现金流或营运资本桥。

\section{利润质量是高空间案例的共同拦截器}

九安医疗的预告驱动来自被投科技企业估值变动，并提示退出与流动性不确定性；吉林敖东和中山公用的预告或年报中投资收益影响较大；辽宁成大亦需要区分联营投资收益与经营现金。此类利润可影响净资产和股价，却不足以在没有可复建资产价值模型时自动形成当前目标价。
